\chapter{Models, Tools and Agents}
\label{ch:declarations}

\section{Models}

\begin{algo}
model smart = "claude-sonnet-5-5"
model fast  = "claude-haiku-5-5"
\end{algo}

A \kw{model} declaration binds a short name to a model identifier. Agents select a model by short name with the \kw{model} clause; an unknown name is error \errcode{E106}. An agent without a \kw{model} clause uses \code{claude-sonnet-5-5}.

\begin{tip}
Name models by their role in your system---\code{smart}, \code{fast}, \code{cheap}---not by product. When a new model is released you change one line, and every agent that uses that role moves with it.
\end{tip}

\section{Tools}

\begin{algo}
tool web.search
tool file.read
\end{algo}

A \kw{tool} declaration announces a capability that the program may grant. Declaring a tool gives it to no one: each agent must list the tools it needs in a \kw{uses} clause, and listing a tool that was never declared is error \errcode{E104}.

\begin{ethics}[Least privilege]
An agent can only be granted capabilities that the program declares openly at the top of the file, and only the agents that list a tool receive it. A reviewer can see every capability of a program, and which agent holds it, without reading any agent's code. This is the security principle of least privilege applied to AI agents.
\end{ethics}

\begin{caution}
Tools are \emph{capability declarations} that the compiler enforces. Since version 0.2 they are also executable: the runtime calls the registered Python function (built-ins \code{calc.eval}, \code{file.read}, \code{kb.search}, or your own plugin) and records every call in the audit log (Chapter~\ref{ch:enterprise}). MCP bindings are planned for version 0.3.
\end{caution}

\section{Agents}

An agent declaration has a name, typed parameters, a return type and a body of clauses:

\begin{algo}[title=agent anatomy]
agent CriterionReviewer(criterion: text, paper: text, depth: text = "detailed") -> text {
    role   "Domain reviewer"
    goal   "Give a {depth} review of {paper} strictly on: {criterion}"
    uses   file.read, web.search
    model  fast
    policy {
        budget tokens 4000
        retry 1
    }
}
\end{algo}

\subsection{Parameters}

Each parameter is \code{name: type}, optionally with a default literal: \code{tone: text = "academic"}. A parameter with a default may be omitted by the caller. Parameter names must not appear in the agent's \kw{forbid} list (\errcode{E108}).

\subsection{Clauses}

\begin{center}
\small
\begin{tabularx}{\linewidth}{@{}L{1.6cm}L{1.8cm}Y@{}}
\toprule
\thead{Clause} & \thead{Required} & \thead{Meaning}\\
\goldrule
\kw{goal} & \textbf{yes} & The work the agent must do, in plain language. Placeholders \code{\{param\}} are filled when it runs. Missing goal: \errcode{E112}.\\
\kw{role} & no & The capacity in which the agent acts; becomes its persona (system prompt in live mode).\\
\kw{uses} & no & Comma-separated tools the agent may use; each must be declared (\errcode{E104}).\\
\kw{model} & no & Short name of a declared model (\errcode{E106}).\\
\kw{policy} & no & Governance rules---see Chapter~\ref{ch:policies}.\\
\bottomrule
\end{tabularx}
\end{center}

Clauses may appear in any order. A clause may be repeated; for \kw{uses} the lists are combined, and for the others the last one wins.

\subsection{What the model receives}

At run time the runtime renders the agent into a prompt with three parts:

\begin{plaincode}[title=rendered prompt]
ROLE: Domain reviewer
GOAL: Give a detailed review of LG-CEF … strictly on: Methodology & rigour
TOOLS: file.read, web.search
\end{plaincode}

In live mode the role becomes the system persona and the model is asked to reply with a JSON object holding its \code{value} and \code{confidence}, in the shape required by the agent's return type.

\section{Writing good agents}

\begin{enumerate}
  \item \textbf{One job per agent.} If a goal contains ``and then'', split it into two agents and connect them in a flow.
  \item \textbf{Write goals as instructions to a capable colleague.} State the output you expect: length, format, audience.
  \item \textbf{Choose the narrowest return type.} If the flow will branch on it, return a \kw{score} or \kw{bool}, not \kw{text}.
  \item \textbf{Grant the fewest tools.} An agent that only rewrites text needs no tools at all.
  \item \textbf{Always set a budget} for agents that will run live.
  \item \textbf{Gate decisions that affect people} with \kw{require human when confidence < …}.
\end{enumerate}
